\documentclass[11pt,a4paper]{article}
\usepackage{amsmath,amssymb,amsthm}
\newtheorem{theorem}{Theorem}[section]
\begin{document}

\section{Statements}

Inline math like $e^{i\pi}+1=0$ and $\sum_{k=1}^{n} k = \tfrac{n(n+1)}{2}$
inside running text, followed by displayed equations:
\begin{equation}\label{eq:gauss}
  \int_{-\infty}^{\infty} e^{-x^2}\,dx = \sqrt{\pi}
\end{equation}

\begin{theorem}\label{thm:main}
  For every $\varepsilon>0$ there exists $\delta>0$ such that
  $|f(x)-f(a)|<\varepsilon$ whenever $|x-a|<\delta$.
\end{theorem}

\begin{align}
  a_1 &= b_1 + c_1, \\
  a_2 &= b_2 + c_2 + d_2, \label{eq:second}\\
  a_3 &= b_3.
\end{align}

\section{Uses}

Equation~\eqref{eq:gauss} and~\eqref{eq:second} together with
Theorem~\ref{thm:main} give the usual argument, spelled out here in prose
long enough to form a real paragraph around the mathematics.

Matrices too:
\[
  M = \begin{pmatrix} 1 & 2 \\ 3 & 4 \end{pmatrix},\qquad
  \det M = -2 .
\]

\end{document}
